\documentclass[11pt, a4paper]{article}

\input{bfw-style.tex}

\title{\vspace*{-1.5cm}\Huge\bfseries\color{bfwPrimaryDark}Aufgabenheft Session~7\\[0.4em]
  \large\color{bfwInkSoft}\textnormal{Präsentation vorbereiten \& planen --- Übungen im IHK-Klausurformat}}
\author{\textsf{IT Lernfeld 1 --- Das Unternehmen und die eigene Rolle im Betrieb beschreiben}}
\date{Selbstlernmaterial}

\begin{document}
\maketitle
\thispagestyle{empty}

\begin{merksatz}[title={\bfseries So nutzen Sie dieses Aufgabenheft}]
Bearbeiten Sie alle Aufgaben zunächst \emph{ohne} in die Lösungen zu schauen. Schreiben
Sie Ihre Antworten in vollständigen Sätzen --- die IHK-Prüfung verlangt formulierte
Begründungen, nicht bloße Stichworte. Beachten Sie die \emph{Operatoren} (nennen,
erläutern, beschreiben, begründen, beurteilen) --- sie geben den geforderten
Antwort-Tiefegrad vor. Erst danach öffnen Sie den Lösungsteil ab Seite~\pageref{loesungen}
und vergleichen.
\end{merksatz}

\bigskip

\textbf{Kontext:} Sie sind Auszubildende bei der \textbf{JIKU IT-Solutions GmbH} in Hamburg.
Ihr Ausbildungsleiter Herr Brandt hat das Team beauftragt, den Ausbildungsbetrieb vor einer
Gruppe von Schülerinnen und Schülern zu präsentieren, die einen Betriebsbesuch planen.
Das Team hat 90 Minuten zur Vorbereitung und hält die Präsentation in drei Wochen.

\bigskip

\noindent
\begin{tabular}{|l|l|l|l|l|l|l|l|l|}
\hline
\textbf{Aufgabe} & 1 & 2 & 3 & 4 & 5 & 6 & 7 & \textbf{Gesamt}\\
\hline
\textbf{Punkte} & /8 & /10 & /8 & /10 & /10 & /6 & /10 & /62\\
\hline
\end{tabular}

\tableofcontents
\vspace{1em}
\hrule

\section{Aufgaben}

\subsection*{Aufgabe~1 --- IPERKA-Methode (8 Punkte)}

\begin{enumerate}[label=(\alph*)]
  \item \textbf{Nennen} Sie die sechs Phasen der IPERKA-Methode und \textbf{erläutern} Sie
  jede Phase kurz. \hfill (6~P.)

  \vspace{9cm}

  \item \textbf{Ordnen} Sie die folgenden Arbeitsschritte der richtigen IPERKA-Phase zu. \hfill (2~P.)

  \bigskip

  \noindent
  \begin{tabular}{|p{9cm}|p{3.5cm}|}
  \hline
  \textbf{Arbeitsschritt} & \textbf{IPERKA-Phase}\\
  \hline
  Das Team bespricht nach der Präsentation mit Herrn Brandt, was gut lief und was verbessert werden kann. & \\[0.8cm]
  \hline
  Mia Schulz legt fest, welche Folien bis wann fertig sein müssen und wer welchen Abschnitt übernimmt. & \\[0.8cm]
  \hline
  Leon Kraft hält die Präsentation vor den Besuchern. & \\[0.8cm]
  \hline
  Das Team liest die Aufgabenstellung und fragt Herrn Brandt nach den Anforderungen. & \\[0.8cm]
  \hline
  \end{tabular}
\end{enumerate}

\subsection*{Aufgabe~2 --- Projektteam und Teamrollen (10 Punkte)}

\begin{enumerate}[label=(\alph*)]
  \item \textbf{Beschreiben} Sie, was ein Projektteam von einer einfachen Arbeitsgruppe unterscheidet. \hfill (3~P.)

  \vspace{5cm}

  \item \textbf{Nennen} Sie die fünf Erfolgsfaktoren für ein gutes Projektteam. \hfill (5~P.)

  \vspace{7cm}

  \item Das JIKU-Präsentationsteam hat neun Mitglieder. \textbf{Beurteilen} Sie, ob diese
  Größe optimal ist, und \textbf{begründen} Sie Ihre Antwort. \hfill (2~P.)

  \vspace{4cm}
\end{enumerate}

\subsection*{Aufgabe~3 --- Teamentwicklung nach Tuckman (8 Punkte)}

\begin{enumerate}[label=(\alph*)]
  \item \textbf{Nennen} Sie die fünf Phasen des Teamentwicklungsmodells von Bruce Tuckman
  und ordnen Sie jeweils ein typisches Merkmal zu. \hfill (5~P.)

  \bigskip

  \noindent
  \begin{tabular}{|p{3.5cm}|p{8.5cm}|}
  \hline
  \textbf{Phase (Name)} & \textbf{Typisches Merkmal}\\
  \hline
   & \\[1.0cm]
  \hline
   & \\[1.0cm]
  \hline
   & \\[1.0cm]
  \hline
   & \\[1.0cm]
  \hline
   & \\[1.0cm]
  \hline
  \end{tabular}

  \bigskip

  \item Im JIKU-Team gibt es Streit: Zwei Mitglieder wollen nicht zusammenarbeiten, Beschlüsse
  werden angezweifelt. \textbf{Benennen} Sie die Tuckman-Phase und \textbf{erläutern} Sie,
  wie Projektleiterin Ana Kovac darauf reagieren sollte. \hfill (3~P.)

  \vspace{6cm}
\end{enumerate}

\subsection*{Aufgabe~4 --- Zielgruppenanalyse und Rahmenbedingungen (10 Punkte)}

\textit{Die Präsentation findet im JIKU-Schulungsraum statt: Beamer vorhanden, 45 Minuten Zeit
inklusive Fragen, kein Druckkostenbudget. Die Besucher sind Schüler der 10.~Klasse
(ca. 20 Personen, kaum IT-Kenntnisse, Interesse an Ausbildungsberufen).}

\begin{enumerate}[label=(\alph*)]
  \item \textbf{Beschreiben} Sie die Zielgruppe (Vorkenntnisse, Erwartungen, Interessen) und
  \textbf{leiten} Sie daraus zwei konkrete Konsequenzen für Inhalt und Sprache ab. \hfill (4~P.)

  \vspace{6cm}

  \item Erstellen Sie eine Übersicht der Rahmenbedingungen für die beschriebene Situation. \hfill (4~P.)

  \bigskip

  \noindent
  \begin{tabular}{|p{3.5cm}|p{9cm}|}
  \hline
  \textbf{Faktor} & \textbf{Situation bei JIKU}\\
  \hline
  Kosten & \\[0.9cm]
  \hline
  Zeit & \\[0.9cm]
  \hline
  Vorkenntnisse & \\[0.9cm]
  \hline
  Raum/Medien & \\[0.9cm]
  \hline
  Unterlagen & \\[0.9cm]
  \hline
  \end{tabular}

  \bigskip

  \item Wie viele Folien sind laut 10-20-30-Regel maximal zulässig? \textbf{Erläutern} Sie,
  was die Regel außerdem vorschreibt. \hfill (2~P.)

  \vspace{3.5cm}
\end{enumerate}

\subsection*{Aufgabe~5 --- Gliederung und roter Faden (10 Punkte)}

\begin{enumerate}[label=(\alph*)]
  \item \textbf{Beschreiben} Sie die drei Strukturabschnitte einer Präsentation
  (Einstieg, Hauptteil, Schluss) und nennen Sie je eine wichtige Funktion. \hfill (6~P.)

  \vspace{8cm}

  \item Das JIKU-Team plant folgende Folienreihenfolge:\\[0.4em]
  \textit{Folie 1: Produktportfolio · Folie 2: Firmenname und Logo · Folie 3: Ausbildungsberufe
  · Folie 4: Teamvorstellung · Folie 5: Außenansicht des Gebäudes}

  \textbf{Beurteilen} Sie, ob diese Reihenfolge einen roten Faden ergibt, und
  \textbf{schlagen} Sie eine verbesserte Gliederung vor. \hfill (4~P.)

  \vspace{6cm}
\end{enumerate}

\subsection*{Aufgabe~6 --- Medienwahl (6 Punkte)}

\textbf{Nennen} Sie drei unterschiedliche Visualisierungsmedien und \textbf{beschreiben} Sie
je einen Vorteil und einen Nachteil in Bezug auf die JIKU-Betriebspräsentation. \hfill (6~P.)

\bigskip

\noindent
\begin{tabular}{|p{3cm}|p{5.5cm}|p{4cm}|}
\hline
\textbf{Medium} & \textbf{Vorteil} & \textbf{Nachteil}\\
\hline
 & & \\[1.2cm]
\hline
 & & \\[1.2cm]
\hline
 & & \\[1.2cm]
\hline
\end{tabular}

\subsection*{Aufgabe~7 --- Fallaufgabe: Betriebspräsentation planen (10 Punkte)}

Sie erstellen ein \textbf{Storyboard} (Folienskizze) für eine fünfminütige Präsentation der
JIKU IT-Solutions GmbH vor Schülerinnen und Schülern.

\begin{enumerate}[label=(\alph*)]
  \item \textbf{Skizzieren} Sie fünf Folien: Notieren Sie zu jeder Folie einen Titel und
  maximal drei Stichpunkte. Beachten Sie Einstieg --- Hauptteil --- Schluss. \hfill (6~P.)

  \bigskip

  \noindent
  \begin{tabular}{|p{0.5cm}|p{3.5cm}|p{8.5cm}|}
  \hline
  \textbf{Nr.} & \textbf{Folientitel} & \textbf{Stichpunkte (max. 3)}\\
  \hline
  1 & & \\[1.5cm]
  \hline
  2 & & \\[1.5cm]
  \hline
  3 & & \\[1.5cm]
  \hline
  4 & & \\[1.5cm]
  \hline
  5 & & \\[1.5cm]
  \hline
  \end{tabular}

  \bigskip

  \item \textbf{Wählen} Sie für jede der fünf Folien ein passendes Visualisierungselement
  (Foto, Diagramm, Aufzählung, Tabelle o.\,Ä.) und \textbf{begründen} Sie Ihre Wahl kurz. \hfill (4~P.)

  \vspace{7cm}
\end{enumerate}

% =====================================================================
\newpage
\section*{Lösungshinweise für Lehrende}
\label{loesungen}
\addcontentsline{toc}{section}{Lösungshinweise für Lehrende}
% =====================================================================

\begin{merksatz}[title={\bfseries Hinweis zur Nutzung}]
Dieser Abschnitt ist ausschließlich für Lehrende bestimmt. Die Lösungen geben
Musterlösungen und typische Fehlerquellen an. Teilnehmende erhalten diesen Teil
erst nach der eigenständigen Bearbeitung.
\end{merksatz}

\subsection*{Lösung Aufgabe~1}

\begin{enumerate}[label=(\alph*)]
  \item Die sechs IPERKA-Phasen:
  \begin{itemize}
    \item \textbf{I --- Informieren:} Aufgabe verstehen; benötigte Informationen recherchieren oder erfragen.
    \item \textbf{P --- Planen:} Arbeitsschritte, Zeitbedarf, Hilfsmittel, Teilziele und Termine festhalten; Maßnahmenplan erstellen.
    \item \textbf{E --- Entscheiden:} Festlegen, welcher Plan umgesetzt wird; Alternativen abwägen; Projektleitung bestimmen.
    \item \textbf{R --- Realisieren:} Den Arbeitsplan wie vereinbart umsetzen.
    \item \textbf{K --- Kontrollieren:} Ergebnisse prüfen, bevor sie dem Auftraggeber übergeben werden; Selbst- und Fremdkontrolle anhand eines Kriterienbogens.
    \item \textbf{A --- Auswerten:} Mit dem Auftraggeber reflektieren, was gut lief und wo Verbesserungspotenzial liegt.
  \end{itemize}

  \textit{Typischer Fehler: Das „A" (Auswerten) wird oft vergessen oder mit dem Kontrollieren verwechselt. Auswerten findet \emph{mit dem Auftraggeber} statt.}

  \item Zuordnung:
  \begin{itemize}
    \item Nachbesprechung mit Herrn Brandt: \textbf{A --- Auswerten}
    \item Mia legt Termine und Verantwortungen fest: \textbf{P --- Planen}
    \item Leon hält die Präsentation: \textbf{R --- Realisieren}
    \item Team liest Aufgabenstellung und fragt Herrn Brandt: \textbf{I --- Informieren}
  \end{itemize}
\end{enumerate}

\subsection*{Lösung Aufgabe~2}

\begin{enumerate}[label=(\alph*)]
  \item Bei der \textbf{einfachen Gruppenarbeit} arbeiten alle gleichwertig und wenig arbeitsteilig an einer gemeinsamen Aufgabe. Ein \textbf{Projektteam} setzt sich aus Personen mit \emph{unterschiedlichen Kompetenzen} zusammen, arbeitet auf ein befristetes besonderes Ziel hin und benötigt eine Projektleitung, die koordiniert und den Gesamtüberblick behält.

  \item Die fünf Erfolgsfaktoren:
  \begin{enumerate}
    \item Klares Ziel
    \item Richtige Teamgröße (5--8 Personen)
    \item Unterschiedliche Teampersönlichkeiten
    \item Akzeptierte Projektleitung
    \item Funktionierende Kommunikation
  \end{enumerate}

  \item Neun Mitglieder liegen leicht über der optimalen Größe von 5--8 Personen. Es entstehen mehr Reibungsverluste und Einigungsprozesse dauern länger. Empfehlung: Das Team in zwei Teilgruppen (z.\,B.\ Planung und Gestaltung) aufteilen, die jeweils eigenverantwortlich arbeiten.
\end{enumerate}

\subsection*{Lösung Aufgabe~3}

\begin{enumerate}[label=(\alph*)]
  \item Die fünf Phasen nach Tuckman:
  \begin{itemize}
    \item \textbf{Forming} (Einstieg): Kennenlernen, höflicher Umgang, erste Regeln.
    \item \textbf{Storming} (Streit): Unterschwellige Konflikte, Cliquenbildung, Sacharbeit tritt in den Hintergrund.
    \item \textbf{Norming} (Regelung): Regeln werden diskutiert und anerkannt; Rollen festigen sich.
    \item \textbf{Performing} (Leistung): Zielorientiert, ideen­reich, produktiv; gegenseitige Wertschätzung.
    \item \textbf{Adjourning} (Auflösung): Abschluss des Projekts, Reflexion der Teamarbeit, Würdigung der gemeinsamen Leistung.
  \end{itemize}

  \item Die Situation beschreibt die \textbf{Storming-Phase}. Ana Kovac sollte: (1) Konflikt offen ansprechen, ohne zu eskalieren; (2) Klare Regeln vereinbaren (z.\,B.\ Entscheidungen durch Mehrheit); (3) Rollen klar zuweisen mit namentlicher Verantwortung; (4) Eine kurze Feedbackrunde einplanen, um Spannungen konstruktiv abzubauen. Die Phase ist normal und geht bei guter Moderation in die Norming-Phase über.
\end{enumerate}

\subsection*{Lösung Aufgabe~4}

\begin{enumerate}[label=(\alph*)]
  \item Zielgruppe: Schüler der 10.~Klasse, kaum IT-Kenntnisse, interessiert an Ausbildungsberufen. Konsequenzen: (1)~\textbf{Einfache Sprache} ohne Fachbegriffe oder mit Erklärung bei jedem Fachbegriff --- keine Annahme von IT-Vorwissen. (2)~\textbf{Inhaltliche Breite} statt Tiefe: Überblick über Betrieb, Ausbildungsberufe und Einstiegsmöglichkeiten statt technischer Details.

  \item Rahmenbedingungen bei JIKU:
  \begin{itemize}
    \item \textbf{Kosten:} Kein Druckkostenbudget --- kein Handout in Papierform; ggf.\ digitale Version.
    \item \textbf{Zeit:} 45 Minuten gesamt inklusive Fragen --- ca. 30--35 Minuten reiner Vortrag.
    \item \textbf{Vorkenntnisse:} Kaum IT-Kenntnisse; allgemeinverständliche Sprache erforderlich.
    \item \textbf{Raum/Medien:} Schulungsraum mit Beamer --- Folienpräsentation möglich.
    \item \textbf{Unterlagen:} Keine gedruckten Unterlagen; ggf.\ Link zu digitalen Materialien.
  \end{itemize}

  \item Laut 10-20-30-Regel: maximal \textbf{10 Folien}. Die Regel schreibt außerdem vor: Vortragsdauer max.\ 20 Minuten und keine Schrift kleiner als 30 Punkt.
\end{enumerate}

\subsection*{Lösung Aufgabe~5}

\begin{enumerate}[label=(\alph*)]
  \item Die drei Strukturabschnitte:
  \begin{itemize}
    \item \textbf{Einstieg:} Interesse wecken, Thema nennen, Überblick geben. Die ersten Sätze prägen sich dem Publikum am stärksten ein.
    \item \textbf{Hauptteil:} Sachlogisch gegliederte Kernaussagen; Spannungsbogen halten; Beispiele verwenden. Zeitanteil ca. 70--75\,\% (praxisübliche Faustregel).
    \item \textbf{Schluss:} Zusammenfassung, Fazit, Ausblick. Die letzten Sätze bleiben ebenfalls lang in Erinnerung.
  \end{itemize}

  \item Die vorliegende Reihenfolge ergibt \textbf{keinen roten Faden}: Die Präsentation beginnt mit Details (Produktportfolio) statt mit der Unternehmensvorstellung. Verbesserte Gliederung: Folie~1: Firmenname, Standort und Logo (Einstieg/Überblick) → Folie~2: Außen-/Innenansicht (Orientierung) → Folie~3: Branche und Hauptgeschäftsfelder (Kerngeschäft) → Folie~4: Ausbildungsberufe bei JIKU (Hauptinteresse der Zielgruppe) → Folie~5: Teamvorstellung und Kontakt (Abschluss).
\end{enumerate}

\subsection*{Lösung Aufgabe~6}

\begin{itemize}
  \item \textbf{Beamer-Präsentation:} Vorteil: Strukturiert, gut sichtbar für alle, Fotos und Diagramme möglich. Nachteil: Technikabhängigkeit; bei Ausfall kein Ersatz.
  \item \textbf{Flip-Chart:} Vorteil: Flexibel, spontane Ergänzungen möglich, kein Technikrisiko. Nachteil: Für mehr als 15 Personen schwer lesbar; Erstellung ist zeitaufwendig.
  \item \textbf{Handout:} Vorteil: Verbleibt bei den Teilnehmern; ermöglicht Nacharbeitung. Nachteil: Kein Druckkostenbudget vorhanden --- bei dieser JIKU-Situation daher nicht einsetzbar.
\end{itemize}

\subsection*{Lösung Aufgabe~7}

\begin{enumerate}[label=(\alph*)]
  \item Mögliches Storyboard:
  \begin{enumerate}
    \item \textbf{Einstieg: Wir sind JIKU IT-Solutions} --- Firmenname, Standort Hamburg, Gründungsjahr
    \item \textbf{Was wir tun: IT-Dienstleistungen und -Produkte} --- Softwareentwicklung, IT-Beratung, Managed Services
    \item \textbf{Unsere Ausbildungsberufe} --- Fachinformatiker/-in (Systemintegration / Anwendungsentwicklung), Kaufmann/-frau für IT-System-Management
    \item \textbf{Ein Tag als Azubi bei JIKU} --- Typische Aufgaben, Team, Lernmöglichkeiten
    \item \textbf{Zusammenfassung und Kontakt} --- Kernbotschaft, nächste Schritte (Bewerbung), Ansprechperson
  \end{enumerate}

  \item Visualisierungselemente:
  Folie 1: \textbf{Logo} --- Wiedererkennungswert; Folie 2: \textbf{Tabelle} (Leistungsübersicht) --- strukturierter Vergleich; Folie 3: \textbf{Aufzählung mit Icons} --- schnelle Erfassbarkeit; Folie 4: \textbf{Foto} (Azubis im Büro) --- persönliche Nähe, Authentizität; Folie 5: \textbf{QR-Code} (Bewerbungsseite) --- direkter Handlungsaufruf.
\end{enumerate}

\end{document}
